% from ega4-1-fr.tex
\oldpage[IV-1]{249}

\subsection{Application aux morphismes ouverts}
\label{subsection:IV.1.10-fr}

\begin{theorem}[1.10.1]
\label{IV.1.10.1-fr}
Soient $X$ un préschéma, $E$ une partie ind-constructible de $X$ (1.9.4), $x$
un point de $X$. Pour que $x$ soit intérieur à $E$, il faut et il suffit que
toute générisation ($0_{\mathrm I}$, 2.1.2) $x'$ de $x$ appartienne à $E$,
autrement dit (I, 2.4.2) que l'on ait
$\operatorname{Spec}(\mathcal O_{X,x})\subset E$.
\end{theorem}

\begin{proof}
La condition est évidemment nécessaire, tout voisinage de $x$ contenant les
générisations de $x$. Pour voir qu'elle est suffisante, on peut évidemment se
borner au cas où $X=\operatorname{Spec}(A)$ est affine, $x$ étant un idéal premier
$\mathfrak p$ de $A$. On sait alors (1.9.5, (ix)) qu'il existe un schéma affine
$X'=\operatorname{Spec}(A')$ et un morphisme $f:X'\to X$ tels que
$f(X')=X-E$. Posons $Y=\operatorname{Spec}(\mathcal O_{X,x})$ et
$Y'=X'\times_XY$~; l'hypothèse signifie (en
vertu de (I, 3.6.5)) que l'on a $Y'=\varnothing$, autrement dit
$A'\otimes_AA_{\mathfrak p}=0$. Comme
$A_{\mathfrak p}=\varinjlim A_t$, où $t$ parcourt $A-\mathfrak p$
($0_{\mathrm I}$, 1.4.5), on a $\varinjlim A_t'=0$, le foncteur
$\varinjlim$ commutant au produit tensoriel. Par suite (1.9.1), il existe un
$t\in A-\mathfrak p$ tel que $A_t'=0$, et $\mathrm D(t)$ est alors un
voisinage ouvert de $x$ dans $X$ contenu dans $E$.
\end{proof}

\begin{corollary}[1.10.2]
\label{IV.1.10.2-fr}
Soient $f:X\to Y$ un morphisme localement de présentation finie, et soit
$y\in Y$. Pour que $y$ soit intérieur à $f(X)$ il faut et il suffit que toute
générisation $y'$ de $y$ appartienne à $f(X)$.
\end{corollary}

\begin{proof}
On sait en effet (1.9.5, (viii)) que $f(X)$ est ind-constructible dans $Y$ et
il suffit d'appliquer (1.10.1).
\end{proof}

Nous dirons qu'une application continue $f:X\to Y$ est \emph{ouverte en un
point} $x\in X$ si l'image par $f$ de \emph{tout} voisinage de $x$ dans $X$
est un voisinage de $f(x)$ dans $Y$.

\begin{proposition}[1.10.3]
\EGAReaderTarget{ega4.c04.prop.1-10-3}{IV.1.10.3-fr}
Soient $f:X\to Y$ un morphisme, $x$ un point de $X$, $y=f(x)$. Considérons
les conditions suivantes~:
\begin{enumerate}
  \item[a)] $f$ est ouvert au point $x$.
  \item[b)] Pour toute générisation $y'$ de $y$, il existe une générisation
  $x'$ de $x$ telle que $f(x')=y'$.
  \item[b')] L'image par $f$ de $\operatorname{Spec}(\mathcal O_{X,x})$ est
  $\operatorname{Spec}(\mathcal O_{Y,y})$.
  \item[c)] Pour toute partie fermée irréductible $Y'$ de $Y$ contenant $y$,
  il existe une composante irréductible de $X'=f^{-1}(Y')$ contenant $x$ et
  dominant $Y'$.
\end{enumerate}
Alors on a les implications
$\mathrm{a)}\Rightarrow\mathrm{b)}\Leftrightarrow\mathrm{b')}
\Leftrightarrow\mathrm{c)}$. Si en outre $f$ est localement de présentation
finie, les quatre conditions sont équivalentes.
\end{proposition}

\begin{proof}
Par définition d'une générisation, l'image par $f$ de
$\operatorname{Spec}(\mathcal O_{X,x})$ est toujours contenue dans
$\operatorname{Spec}(\mathcal O_{Y,y})$, donc b) et b') sont équivalentes. Il
est immédiat
que c) entraîne b), car si $y'$ est une générisation de $y$,
$Y'=\overline{\{y'\}}$ est une partie irréductible de $Y$ contenant $y$~; si
$x'$ est le point générique d'une composante irréductible de $f^{-1}(Y')$ qui
contient $x$ et domine $Y'$, on a $f(x')=y'$ ($0_{\mathrm I}$, 2.1.5) et $x'$
est une générisation de $x$. Inversement, b) entraîne c)~: en effet, soit
$y'$ le point générique de $Y'$ et soit $x'$ une générisation de $x$ telle
que $f(x')=y'$~; soit $Z'$ une composante irréductible de $f^{-1}(Y')$
contenant $x'$ (donc $x$)~; comme $f(x')$ est déjà dense dans $Y'$, il en est
ainsi \emph{a fortiori} de $f(Z')$.

Pour voir que a) entraîne b'), on peut évidemment se borner au cas où
$X=\operatorname{Spec}(B)$ et $Y=\operatorname{Spec}(A)$ sont affines, $x$
étant un idéal premier
$\mathfrak q$ de $B$, de sorte que
$\mathcal O_{X,x}=B_{\mathfrak q}=\varinjlim B_t$,
